\documentclass[25pt, a0paper, portrait]{tikzposter}
\title{\parbox{\linewidth}{\centering Multi-Port Memory with Bidirectional Ports for FPGAs Using XOR and LVT Methods}}
  \author{Kevin Townsend}
  \date{Feb. 23rd, 2026}
\usetheme{Board}

\usepackage{tikz}
\usepackage[style=numeric,backend=bibtex]{biblatex}
\usepackage{fp}
\usepackage{ifthen}
\usepackage{makecell}
\usepackage{booktabs}
\addbibresource{poster.bib}


\begin{document}

\maketitle[width=\paperwidth]

\begin{columns}
  \column{0.67}
  %TODO: Motivation
  \block{Motivation}{
As computation needs keep increasing, one way to keep up has been
specialized architectures. FPGAs provide a way to implement such architectures without taping out an ASIC.
However, the limitations of FPGA resources requires some creativity to map designs to FPGAs.
This paper explores how to overcome the limitation that FPGA memories have a limited number of
ports. Specifically, we propose a method to create memories with more than 2 ports.
  }
  \block{Source Code}{
We provide all of the source code used in implementation and testing our design at
https://github.com/kevintownsend/mpm. We tested our design with Verilator\cite{verilator_manual} and synthesized the design
with Vivado\cite{ug893} targeting an AMD Virtex UltraScale+ HBM VU47P-3 FPGA.

  }
  %TODO: Generalizing XOR memory
  \block{Generalizing XOR Memory}{

  \begin{tikzpicture}[scale=3.8]
  \foreach \x in {0,1,2,3,4,5} {
    \foreach \y in {0,1,2,3} {
      \ifthenelse{\(\x = 2 \AND \y = 1\) \OR \(\x = 3 \AND \y = 0\)}{}{
        \node (r\x\y)[trapezium, trapezium left angle=70, trapezium right angle=110, draw] at (2*\x, \y){RAM$_{\x,\fpeval{3-\y}}$};
      }
    }
  }
  \foreach \x in {0,1,4,5} {
    \node (bport\x)[] at (2*\x+1, -1){PORT\x};
  }
  \foreach \y in {0,1,2,3} {
    \node (lport\y)[] at (-1.3, 3.5-\y){PORT\y};
  }
  \node [rotate=90] at (-2, .8){WRITE PORTS};
  \node [rotate=90] at (-2, 2.8){FULL PORTS};
  \node [] at (1.5, -1.5){FULL PORTS};
  \node [] at (9.5, -1.5){READ PORTS};
  \node(x50)[inner sep=0pt] at (10.8, 0){$\oplus$};
  \node(x51)[inner sep=0pt] at (10.8, 1){$\oplus$};
  \node(x52)[inner sep=0pt] at (10.8, 2){$\oplus$};
  \draw[->] (r50) -- (x50);
  \draw[->] (r51) -- (x51);
  \draw[->] (r52) -- (x52);
  \draw[->] (r53) -| (x52);
  \draw[->] (x52) -- (x51);
  \draw[->] (x51) -- (x50);
  \draw[->] (x50) -- (bport5.130);
  \node(x40)[inner sep=0pt] at (8.8, 0){$\oplus$};
  \node(x41)[inner sep=0pt] at (8.8, 1){$\oplus$};
  \node(x42)[inner sep=0pt] at (8.8, 2){$\oplus$};
  \draw[->] (r40) -- (x40);
  \draw[->] (r41) -- (x41);
  \draw[->] (r42) -- (x42);
  \draw[->] (r43) -| (x42);
  \draw[->] (x42) -- (x41);
  \draw[->] (x41) -- (x40);
  \draw[->] (x40) -- (bport4.130);
  \node(x31)[inner sep=0pt] at (6.8, 1){$\oplus$};
  \node(x32)[inner sep=0pt] at (6.8, 2){$\oplus$};
  \node(x30)[inner sep=0pt] at (6.8, .5){$\oplus$};
  \draw[->] (r31) -- (x31);
  \draw[->] (r32) -- (x32);
  \draw[->] (r33) -| (x32);
  \draw[->] (x32) -- (x31);
  \draw[->] (x31) -- (x30);
  %PORT3 connections.
  \draw[->] (lport3) -- (x30);
  \draw[->] (x30) |- (-.6,.3) |- (r00);
  \draw[->] (1.4,.3) |- (r10);
  \draw[->] (3.4,.3) |- (r20);
  \draw[->] (6.8, .3) -| (7.4,0) -- (r40);
  \draw[->] (7.4, .3) -| (9.4,0) -- (r50);
  %PORT2 XORs
  \node(x22)[inner sep=0pt] at (4.8, 2){$\oplus$};
  \node(x21)[inner sep=0pt] at (4.8, 1.5){$\oplus$};
  %PORT2 connections
  \draw[->] (lport2) -- (x21);
  \draw[->] (r23) -| (x22);
  \draw[->] (r22) -- (x22);
  \draw[->] (x22) -- (x21);
  \draw[->] (r20) -| (x21);
  \draw[->] (x21) -| (5,1.3) -| (-.6,1) -- (r01);
  \draw[->] (1.4,1.3) |- (r11);
  \draw[->] (5,1.3) -- (9.4, 1.3) |- (r51);
  \draw[->] (7.4, 1.3) |- (r41);
  \draw[->] (5.4, 1.3) |- (r31);
  %draw xor of full ports
  %PORT1
  \node(x11)[inner sep=0pt] at (2.8, 1){$\oplus$};
  \node(x12)[inner sep=0pt] at (3.3, 2.5){$\oplus$};
  \node(x12b)[inner sep=0pt] at (3, 2){$\oplus$};
  \node(x13)[inner sep=0pt] at (2.8, 3){$\oplus$};
  \draw[->] (r10) -| (x11);
  \draw[->] (r11) -- (x11);
  \draw[->] (lport1) -- (x12);
  \draw[->] (x11) -- (x13);
  \draw[->] (r13) -- (x13);
  \draw[->] (x13) -| (x12);
  \draw[->] (r12) -- (x12b);
  \draw[->] (x12) -| (3.3, 2.3) -| (-.6,2) -- (r02);
  \draw[->] (1.4,2.3) |- (r12);
  \draw[->] (3.3, 2.3) -- (9.4, 2.3) |- (r52);
  \draw[->] (3.4, 2.3) |- (r22);
  \draw[->] (5.4, 2.3) |- (r32);
  \draw[->] (7.4, 2.3) |- (r42);
  \draw[->] (x13) -| (x12b);
  \draw[->] (x12b) -- (bport1);
  %PORT0
  \node(x01)[inner sep=0pt] at (0.8, 1){$\oplus$};
  \node(x02)[inner sep=0pt] at (0.8, 2){$\oplus$};
  \node(x03)[inner sep=0pt] at (0.8, 3.5){$\oplus$};
  \node(x03b)[inner sep=0pt] at (1.0, 2.8){$\oplus$};
  \draw[->] (r00) -| (x01);
  \draw[->] (r01) -- (x01);
  \draw[->] (x01) -- (x02);
  \draw[->] (x02) -- (x03);
  \draw[->] (lport0) -- (x03);
  \draw[->] (x03) -| (1, 3.3) -| (-.6,3) -- (r03);
  \draw[->] (1, 3.3) -- (9.4, 3.3) |- (r53);
  \draw[->] (1.4, 3.3) |- (r13);
  \draw[->] (3.4, 3.3) |- (r23);
  \draw[->] (5.4, 3.3) |- (r33);
  \draw[->] (7.4, 3.3) |- (r43);
  \draw[->] (.8, 2.8) -- (x03b);
  \draw[->] (r02) -- (x02);
  \draw[->] (r03) -| (x03b);
  \draw[->] (x03b) -- (bport0);
  \end{tikzpicture}
  %\end{scaletikzpicturetowidth}
  %\end{center}

We propose a simple generalization to XOR memories presented in previous work\cite{laforest14}.
XOR memories work by using the following property:
\begin{equation} \label{eqn:xor}
a \oplus b \oplus b = a
\end{equation}
Where $\oplus$ is the bitwise XOR operator. Equation 1 follows from:
\begin{equation}
a \oplus 0 = a
\end{equation}
and
\begin{equation}
a \oplus a = 0
\end{equation}
Meaning XOR is the inverse operator of itself. Combined with the commutative and associative property of XOR we can generalize this to:
\begin{equation}
b_0 \oplus \dots \oplus b_n \oplus (a \oplus b_0 \dots \oplus b_n) = a
\end{equation}


We use XOR logic to add bidirectional
ports and analyze the performance of distributed memory and block memory
versions of this design.

The number of RAMs ($N_{RAM}$) needed for this memory is:
\begin{equation}
N_{RAM}=(W+F)(W+F+R)-W
\end{equation}
Where $W$ is the number of write ports, $R$ is the number of read ports and $F$ is the number of full ports. This expands to:
\begin{equation}
N_{RAM}=W^2 + 2WF + F^2 + WR + FR - W
\end{equation}

    %TODO: complete.
    \begin{minipage}{0.48\linewidth}
        %\begin{tikzfigure}[Caption for Left Figure]
    \centering
    \begin{tikzpicture}[scale=9]
      \draw (.25,1.75) rectangle (.75,2.25);
      \foreach \j in {-.15,-.05, .05, .15}{
          \draw[] (.25, 2+\j) -- (.75,2+\j);
      }
      \foreach \i/\v in {0/1101,1/1111,2/0111,3/1000,4/0101}{
          \FPeval{\x}{.7-\i*.1};
          \FPeval{\y}{2-.2+\i*.1};
          \node[] at (.5,\y) {\v};
      }
      \node at (.5,2.25)[anchor=south]{RAM$_{emulated}$};
      \draw[xshift=-.5cm] (.25,0.75) rectangle (.75,1.25);
      \node at (0,1.25) [anchor=south]{RAM$_{0,0}$};
      \foreach \j in {-.15,-.05, .05, .15}{
          \draw[] (-.25, 1+\j) -- (.25,1+\j);
      }
      \foreach \i/\v in {0/1010,1/1111,2/0100,3/1010,4/0100}{
          \FPeval{\x}{.7-\i*.1};
          \FPeval{\y}{1-.2+\i*.1};
          \node[] at (0,\y) {\v};
      }
      \draw[xshift=.5cm] (.25,0.75) rectangle (.75,1.25);
      \node at (1,1.25) [anchor=south]{RAM$_{1,0}$};
      \foreach \j in {-.15,-.05, .05, .15}{
          \draw[] (1-.25, 1+\j) -- (1.25,1+\j);
      }
      \foreach \i/\v in {0/1010,1/1111,2/0100,3/1010,4/0100}{
          \FPeval{\y}{1-.2+\i*.1};
          \node[] at (1,\y) {\v};
      }
      \draw[xshift=-.5cm,yshift=-1cm] (.25,0.75) rectangle (.75,1.25);
      \node at (0,.25) [anchor=south]{RAM$_{0,1}$};
      \foreach \j in {-.15,-.05, .05, .15}{
          \draw[] (0-.25, 0+\j) -- (0.25,0+\j);
      }
      \foreach \i/\v in {0/0111,1/0000,2/0011,3/0010,4/0001}{
          \FPeval{\y}{1-.2+\i*.1};
          \node[] at (0,-1+\y) {\v};
      }
      \draw[xshift=.5cm,yshift=-1cm] (.25,0.75) rectangle (.75,1.25);
      \node at (1,.25) [anchor=south]{RAM$_{1,1}$};
      \foreach \j in {-.15,-.05, .05, .15}{
          \draw[] (1-.25, 0+\j) -- (1.25,0+\j);
      }
      \foreach \i/\v in {0/0111,1/0000,2/0011,3/0010,4/0001}{
          \FPeval{\y}{1-.2+\i*.1};
          \node[] at (1,-1+\y) {\v};
      }
      %draw XORs
      \node(x0)[inner sep=0pt] at (.4,1.5){$\oplus$};
      \node(x1)[inner sep=0pt] at (1.5,.5){$\oplus$};
      %draw ports
      \node(p0) at (-.5,1.5) [rotate=90]{PORT0};
      \node(p1) at (-.5,.5) [rotate=90]{PORT1};
      %draw connections
      \draw[->] (p1) -- (x1) node[pos=.75,fill=white]{1010};
      \draw[->] (1.25,1.1) -| (x1);
      \draw (x1) |- (-.3, .45);
      \draw[->] (-.3, .45) |- (-.25, 0.1);
      \draw[->] (.65, .45) |- (.75, 0.1);
      \draw[->] (p0) -- (x0) node[pos=.5,fill=white]{1000};
      \draw[->] (.25,-.1) -| (x0);
      \draw[->] (x0) -| (.5,1.45) -| (-.3, .9) -- (-.25, .9);
      \draw[->] (.5,1.45) -| (.7, .9) -- (.75, .9);
      % draw emulated ram connections.
      \node(ep1) at (-.5,2.1) {PORT1};
      \draw[->] (ep1) -- (.25,2.1) node[midway,fill=white]{1010};
      \node(ep0) at (-.5,1.9) {PORT0};
      \draw[->] (ep0) -- (.25,1.9) node[midway,fill=white]{1000};
    \end{tikzpicture}
    \end{minipage}
    \hfill
    \begin{minipage}{0.48\linewidth}
        \centering
    \begin{tikzpicture}[scale=9]
      %draw emulated ram
      %\draw[] (.5,2-.25) -- (1,2-.25) -- (.5,2.25) -- (0, 2.25) -- cycle;
      \draw (.25,1.75) rectangle (.75,2.25);
      \foreach \j in {-.15,-.05, .05, .15}{
          \draw[] (.25, 2+\j) -- (.75,2+\j);
      }
      \foreach \i/\v in {0/1101,1/1000,2/0111,3/1010,4/0101}{
          \FPeval{\x}{.7-\i*.1};
          \FPeval{\y}{2-.2+\i*.1};
          \node[] at (.5,\y) {\v};
      }
      \node at (.5,2.25)[anchor=south]{RAM$_{emulated}$};
      \draw[xshift=-.5cm] (.25,0.75) rectangle (.75,1.25);
      \node at (0,1.25) [anchor=south]{RAM$_{0,0}$};
      \foreach \j in {-.15,-.05, .05, .15}{
          \draw[] (-.25, 1+\j) -- (.25,1+\j);
      }
      \foreach \i/\v in {0/1010,1/1000,2/0100,3/1010,4/0100}{
          \FPeval{\x}{.7-\i*.1};
          \FPeval{\y}{1-.2+\i*.1};
          \node[] at (0,\y) {\v};
      }
      \draw[xshift=.5cm] (.25,0.75) rectangle (.75,1.25);
      \node at (1,1.25) [anchor=south]{RAM$_{1,0}$};
      \foreach \j in {-.15,-.05, .05, .15}{
          \draw[] (1-.25, 1+\j) -- (1.25,1+\j);
      }
      \foreach \i/\v in {0/1010,1/1000,2/0100,3/1010,4/0100}{
          \FPeval{\y}{1-.2+\i*.1};
          \node[] at (1,\y) {\v};
      }
      \draw[xshift=-.5cm,yshift=-1cm] (.25,0.75) rectangle (.75,1.25);
      \node at (0,.25) [anchor=south]{RAM$_{0,1}$};
      \foreach \j in {-.15,-.05, .05, .15}{
          \draw[] (0-.25, 0+\j) -- (0.25,0+\j);
      }
      \foreach \i/\v in {0/0111,1/0000,2/0011,3/0000,4/0001}{
          \FPeval{\y}{1-.2+\i*.1};
          \node[] at (0,-1+\y) {\v};
      }
      \draw[xshift=.5cm,yshift=-1cm] (.25,0.75) rectangle (.75,1.25);
      \node at (1,.25) [anchor=south]{RAM$_{1,1}$};
      \foreach \j in {-.15,-.05, .05, .15}{
          \draw[] (1-.25, 0+\j) -- (1.25,0+\j);
      }
      \foreach \i/\v in {0/0111,1/0000,2/0011,3/0000,4/0001}{
          \FPeval{\y}{1-.2+\i*.1};
          \node[] at (1,-1+\y) {\v};
      }
      %add ports
      \node(p0) at (.5,-.5) {PORT0};
      % add xors
      \node(x0)[inner sep=0pt] at (.5,0.1){$\oplus$};
      %add connections
      \draw[->] (.25, 1.1) -| (x0);
      \draw[->] (.25, .1) -- (x0);
      \draw[->] (x0) -- (p0);
      %add ports
      \node(p1) at (1.5,-.5) {PORT1};
      %add xors
      \node(x1)[inner sep=0pt] at (1.5,-.1){$\oplus$};
      %add connections
      \draw[->] (1.25, .9) -| (x1);
      \draw[->] (1.25, -.1) -- (x1);
      \draw[->] (x1) -- (p1);
      %add emulated ports
      \node(ep0) at (1.25, 2.1) {PORT0};
      \node(ep1) at (1.25, 1.9) {PORT1};
      \draw[->] (.75,2.1) -- (ep0);
      \draw[->] (.75,1.9) -- (ep1);
    \end{tikzpicture}
    \end{minipage}
}
  %TODO: Analysis of XOR Memory Performance
  \column{0.33}
  \block{Analysis of XOR Memory Performance}{
  \begin{tabular}{llllll}
    \toprule
    Ports & LUTs & \thead{LUTs\\ configured\\ as memory} & FFs & BRAMs & Max Frequency\\
    \midrule
    2 & 144 & 0 & 218 & 4 & 625Mhz\\
    4 & 512 & 0 & 444 & 16 & 476Mhz\\
    8 & 1920 & 0 & 920 & 64 & 370Mhz\\
    16 & 6,656 & 0 & 1,968 & 256 & 263Mhz\\
    32 & 10,240 & 0 & 2,400 & 1024 & 123Mhz\\
    \bottomrule
  \end{tabular}
Synthesis results of XOR memory with block RAMs for different port counts. Each memory has a depth of 1024 and width of 32 bits. To reduce the number of IO ports and fit the design on the FPGA, we used a wrapper for the multi-port memory for designs with 16 and 32 ports.
  }
  %TODO: Live Value Table Memory
  \block{Live Value Table Memory}{
  \begin{tikzpicture}[xscale=6, yscale=3]
  %draw rams
  \foreach \x in {0,1,2,3} {
    \foreach \y in {0,1,2,3} {
      \ifthenelse{\x < \y}{
      \node(r\y\x)[trapezium, trapezium left angle=70, trapezium right angle=110, draw] at (\x,3-\y) {RAM$_{\x,\y}$};
      }{
      \node[text=gray,dashed, trapezium, trapezium left angle=70, trapezium right angle=110, draw] at (\x,3-\y) {RAM$_{\x,\y}$};
      }
    }
  }
  %draw a ports
  \foreach \x in {0,1,2,3} {
    %\node(xport\x) at (\x, -1) {\x};
  }
  %draw b ports
  \foreach \y in {0,1,2,3} {
    %\node(yport\y) at (-1, 3-\y) {\y};
  }
  %draw connections
  \draw[<->] (r10.east) -| (.6,-.6) -| (.1,-1);
  \draw[->] (.2,-.6) -- (.2,-1);
  \draw[<->] (r20.east) -| (.5,-.5) -| (.0,-1);
  \draw[<->] (r30.east) -| (.4,-.4) -| (-.1,-1);
  \draw[<->] (r10) |- (1.6, 2.4) |- (1.2, -.6) -- (1.2,-1);
  \draw[->] (1.2,-.6) -| (1.1,-1);
  \draw[<->] (r21) -| (1.5,-.5) -| (1,-1);
  \draw[<->] (r31) -| (1.4,-.4) -| (.9,-1);
  \draw[<->] (r20) |- (2.5,1.5) |- (2.2, -.6) -- (2.2,-1);
  \draw[<->] (r21) |- (2.4,1.4) |- (2.1, -.5) -- (2.1,-1);
  \draw[->] (2.1,-.6) -| (2.0,-1);
  \draw[<->] (r32) -| (2.4,-.4) -| (1.9,-1);
  \draw[<->] (r30) |- (3.2,.6) -- (3.2,-1);
  \draw[<->] (r31) |- (3.1,.5) -- (3.1,-1);
  \draw[<->] (r32) |- (3.0,.4) -- (3.0,-1);
  \draw[->] (3.0,-.6) -| (2.9,-1);
%  %x0 connections
%  \draw (xport0) .. controls (.8,-1) and (.8,2) .. (r10.east);
%  \draw (xport0) .. controls (.6,-1) and (.7,1) .. (r20.east);
%  \draw (xport0) .. controls (.4,-1) and (.6,0) .. (r30.east);
%  %x1 connections
%  \draw (xport1) .. controls (1.6,-1) and (1.7,1) .. (r21);
%  \draw (xport1) .. controls (1.4,-1) and (1.6,0) .. (r31);
%  %x2 connection
%  \draw (xport2) .. controls (2.4,-1) and (2.6,0) .. (r32);
%  %y1 connection
%  \draw (yport1) .. controls (-1, 2.5) and (0, 2.5) .. (r10);
%  %y2 connections
%  \draw (yport2) .. controls (-1, 1.4) and (0, 1.4) .. (r20);
%  \draw (yport2) .. controls (-1, 1.6) and (1, 1.6) .. (r21);
%  %y3 connections
%  \draw (yport3) .. controls (-1, 0.4) and (0, 0.4) .. (r30);
%  \draw (yport3) .. controls (-1, 0.6) and (1, 0.6) .. (r31);
%  \draw (yport3) .. controls (-1, 0.8) and (2, 0.8) .. (r32);
%  %draw labels
%  \node at (-1.5, 1)[rotate=90]{Ports};
%  \node at (1, -1.5){Ports};
  %draw muxes
  \draw[xshift=-.2cm,yshift=-1cm] (0,0) -- (.1,-.2) -- (.4,-.2) -- (.5,0) -- cycle;
  \draw[xshift=.8cm,yshift=-1cm] (0,0) -- (.1,-.2) -- (.4,-.2) -- (.5,0) -- cycle;
  \draw[xshift=1.8cm,yshift=-1cm] (0,0) -- (.1,-.2) -- (.4,-.2) -- (.5,0) -- cycle;
  \draw[xshift=2.8cm,yshift=-1cm] (0,0) -- (.1,-.2) -- (.4,-.2) -- (.5,0) -- cycle;
  %draw LVT
  \node(lvt)[draw] at (1.5, -2) {LVT};
  %draw LVT connections
  \draw[<->] (lvt) -- (.5,-2) |- (.25,-1.1);
  \draw[<->] (lvt.110) |- (1.25,-1.1);
  \draw[<->] (lvt.70) |- (1.85,-1.1);
  \draw[<->] (lvt) -- (2.5,-2) |- (2.85,-1.1);
  %draw ports
  \node[anchor=east] at (.05,-1.4) {PORT};
  \draw[<->] (.05,-1.2) -- (.05,-1.6) node[right,midway]{0};
  \draw[<->] (1.05,-1.2) -- (1.05,-1.6) node[right,midway]{1};
  \draw[<->] (2.05,-1.2) -- (2.05,-1.6) node[right,midway]{2};
  \draw[<->] (3.05,-1.2) -- (3.05,-1.6) node[right,midway]{3};
  \end{tikzpicture}

XOR memories can be used by themselves, but combining a live value table (LVT) with XOR may be more efficient. We present a LVT memory that utilizes XOR memory.

  }
  %TODO: Analysis of LVT Memory
  \block{Analysis of LVT Memory}{
  \begin{tikzpicture}[yscale=2.1, xscale=3]
    \draw[->] (0,0) -- (4.5, 0);
    \draw (0,.1) -- (0, -.1) node[anchor=north] {2};
    \draw (1,.1) -- (1, -.1) node[anchor=north] {4};
    \draw (2,.1) -- (2, -.1) node[anchor=north] {8};
    \draw (3,.1) -- (3, -.1) node[anchor=north] {16};
    \draw (4,.1) -- (4, -.1) node[anchor=north] {32};
    %draw y axis
    \draw[->] (0,0) -- (0, 7.5);
    \draw (1pt, 0) -- (-1pt, 0) node[rotate=45,anchor=south east] {0Mhz};
    \foreach \y in {1,2,3,4,5,6,7}
      \draw (1pt, \y) -- (-1pt, \y) node[rotate=45,anchor=south east] {\y00Mhz};
    %labels
    \node at (2,-1) {PORTS};
    \node at (-1.5,2) [rotate=90]{FREQUENCY};
    %max frequncy
    \draw[dashed] (0, 7.18) -- (4, 7.18) node[anchor=west] {718Mhz};
    %line graphs
    \draw (0, 5.88) -- (1, 5.56) -- (2, 4.55) -- (3, 2.50) -- (4, 1.27);
    \node[fill=blue, draw,label={[rotate=30]right:588Mhz}] at (0,5.88){};
    \node[fill=blue, draw,label={[rotate=30]right:556Mhz}] at (1,5.56){};
    \node[fill=blue, draw,label={[rotate=30]right:455Mhz}] at (2,4.55){};
    \node[fill=blue, draw,label={[rotate=30]right:250Mhz}] at (3,2.50){};
    \node[fill=blue, draw,label={[rotate=30]right:127Mhz}] at (4,1.27){};
    %line graph for pipelnined.
    \draw (0, 7.14) -- (1, 6.67) -- (2, 6.25) -- (3, 4.17) -- (4, 2.47);
    \node[fill=red, draw, circle,inner sep=6pt,label={[rotate=30]right:714Mhz}] at (0,7.14){};
    \node[fill=red, draw, circle,inner sep=6pt,label={[rotate=30]right:667Mhz}] at (1,6.67){};
    \node[fill=red, draw, circle,inner sep=6pt,label={[rotate=30]right:625Mhz}] at (2,6.25){};
    \node[fill=red, draw, circle,inner sep=6pt,label={[rotate=30]right:417Mhz}] at (3,4.17){};
    \node[fill=red, draw, circle,inner sep=6pt,label={[rotate=30]right:247Mhz}] at (4,2.47){};
  \end{tikzpicture}

  Frequency of LVT design. The \protect\tikz{\protect\node[fill=blue, draw]{}}s are the (non-pipelined) LVT memory designs and the \protect\tikz{\protect\node[fill=red, draw, circle,inner sep=6pt]{}}s are the pipelined LVT memory designs.
  }
  %TODO: Bibliography
  \block{Bibliography}{
    \begin{center}\mbox{}\vspace{-\baselineskip}
    \printbibliography
    \end{center}
  }
\end{columns}

\end{document}
